\section{Introduction}
\label{sec:intro}

The integration of large language model (LLM)-backed agents into financial decision loops represents a qualitative shift in the microstructure of depositor behavior. Prior to this development, bank-run dynamics were governed chiefly by human coordination failures under incomplete information, a setting well characterized by the canonical Diamond-Dybvig framework \citep{diamond1983bank}. The collapse of Silicon Valley Bank in March 2023 illustrated how this classical picture is being fundamentally altered: social media amplified panic signals within hours, and depositor action reached a critical threshold far faster than any regulatory circuit-breaker could engage. As financial institutions increasingly deploy AI-assisted portfolio monitoring, cash-flow forecasting, and even automated redemption triggers, the effective decision speed of the depositor population accelerates by orders of magnitude relative to the institutional response capacity. The resulting gap between agent reaction time and stabilization latency constitutes a new and undertheorized systemic risk surface.

Two compounding mechanisms drive this new risk surface. First, the standard Diamond-Dybvig coordination failure is well understood: when a sufficient fraction of depositors anticipate others withdrawing, early exit becomes individually rational regardless of bank solvency, producing a self-fulfilling panic equilibrium. Second, and less studied in the AI context, is the monoculture-induced behavioral correlation that emerges when a large fraction of depositor-agents share an underlying foundation model. Gorton and Pennacchi \citep{gorton1990financial} showed that information-insensitive debt is stable precisely because heterogeneous agents lack the incentive or capacity to coordinate on private information; that insensitivity boundary collapses when agents share a common signal-processing architecture. Acemoglu et al.\ \citep{acemoglu2015systemic} further demonstrated that dense exposure networks dramatically amplify idiosyncratic shocks into systemic cascades. The interaction of these two pre-existing mechanisms in an AI-populated financial system has not been formally analyzed, and it is the gap this paper addresses.

This paper makes three contributions. First, we provide a theoretical integration of the four canonical frameworks---Diamond-Dybvig, Gorton-Pennacchi, Acemoglu, and Shannon---yielding a closed-form lower bound on run probability as a function of fundamental shock magnitude and effective behavioral diversity (Proposition~\ref{prop:run-prob}, Section~\ref{sec:theory}). Second, we construct a synthetic LLM-agent laboratory with a pre-registered experimental protocol (eight experiments, E01--E08) designed to test the proposition's predictions under controlled variation of foundation-model cluster count, correlation structure, time pressure, and prompt framing (Section~\ref{sec:design}). Third, we produce an empirical phase diagram over the $(\sigma_{\text{fund}}, K_{\text{eff}})$ parameter space identifying a critical effective-diversity threshold $K^*$ below which sub-critical fundamentals are sufficient to trigger runs, and quantify run-probability reductions achievable through four counter-mechanism designs (Section~\ref{sec:counter}).

% TODO: tighten Intro to 1.5 pages after results land.
